\begin{table}[htbp]
\centering
\caption{Representative perturbation, counterfactual, and corruption-based evaluation studies.}
\label{tab:perturbation_methods}
\small
\begin{tabularx}{\columnwidth}{@{}p{1.8cm}XX@{}}
\toprule
\textbf{Reference} & \textbf{Method} & \textbf{Application} \\
\midrule
\cite{Zhang2021} & Input perturbation + attribution & ECG diagnosis \\
\cite{MorenoSanchez2023} & Feature ablation & HF survival \\
\cite{singla2023explaining} & Counterfactual generation & Medical imaging \\
\cite{lenis2020domain} & Counterfactual impact & Image interpretation \\
\cite{palatnik2021xai} & Bias-sensitive saliency & COVID CT \\
\cite{Pertuz2023} & Saliency-mask comparison & Breast lesion \\
\cite{Repetto2025EvaluatingTR} & Corruption robustness & Medical imaging \\
\cite{Patricio2024} & AOPC/MoRF perturbation & Image classification \\
\cite{sun2023improving} & Patch perturbation & COVID-19 CXR \\
\cite{jacob2025indian} & Cross-corpus perturbation & Wound healing \\
\cite{akgundougdu2025explainable} & Multi-method perturbation & Brain tumour \\
\cite{alpsalaz2026explainability} & Reliability-weighted perturbation & Cervical cancer \\
\bottomrule
\end{tabularx}
\end{table}
